% Generated draft: requires mathematical and source review.
\documentclass[11pt]{article}
\usepackage[margin=25mm]{geometry}
\usepackage{fontspec}
\setmainfont{Latin Modern Roman}
\usepackage{amsmath,amssymb,mathtools,unicode-math}
\setmathfont{Latin Modern Math}
\usepackage{enumitem,needspace,xurl,hyperref}
\setlength{\parindent}{0pt}
\setlength{\parskip}{5pt}
\setlength{\emergencystretch}{3em}
\newcommand{\N}{\mathbb{N}}
\newcommand{\Z}{\mathbb{Z}}
\newcommand{\Q}{\mathbb{Q}}
\newcommand{\R}{\mathbb{R}}
\newcommand{\C}{\mathbb{C}}
\newcommand{\F}{\mathbb{F}}
\newcommand{\A}{\mathbb{A}}
\newcommand{\PP}{\mathbb{P}}
\newcommand{\E}{\mathbb{E}}
\newcommand{\eps}{\varepsilon}
\newcommand{\dd}{\,\mathrm{d}}
\begin{document}
\section*{UnsolvedMath ID: 6700002}
\textbf{Source problem code:} catalogue entry
\subsection{Definitions and mathematical statement}
For a fixed smooth manifold \(M\), let \(\mathcal R(M)\) be the space of smooth Riemannian metrics, and define
\[
 \mathcal R^+(M)=\{g\in\mathcal R(M):\operatorname{Scal}_g(x)>0
 \text{ for every }x\in M\}.
\]
The source asks for the topology of spaces of this kind. Under the usual smooth-topology reading, \(\mathcal R(M)\) carries the \(C^\infty\) topology and \(\mathcal R^+(M)\) its subspace topology. Determine their topological structure as the underlying manifold varies, including components, homotopy groups based at a chosen \(g_0\in\mathcal R^+(M)\), and, where possible, the full homotopy type. This is a broad question, not a claim that a particular invariant has already been specified as a complete answer.

The excerpt supplies no dimension, compactness, completeness or boundary conditions on \(M\), no normalization such as volume, and no topology on the metric space. On noncompact manifolds compact-open and stronger Whitney-type smooth topologies need not give the same question. The displayed reading concerns the space of metrics itself; no quotient by diffeomorphisms or conformal transformations is inserted. Existence and the possibility that \(\mathcal R^+(M)\) is empty must be kept separate from its nonempty-space topology.

The supplied summary reports substantial existence/surgery and higher-homotopy progress, including stabilized results via what it calls \(\lambda\)-invariants and families index theory, while saying the full homotopy type remains unknown in general, especially in dimension four. It provides no hypotheses or exact theorem citations, so these descriptions are not promoted to a universal existence classification or computation. The 2017 open-item tag and later partial-progress label are source assessments without independent current verification.
\subsection{Short English statement}
Describe the topology of the space of positive-scalar-curvature metrics on a manifold, beyond deciding whether such metrics exist. The underlying manifold class, topology on metrics, and any normalization need specification; the source reports extensive partial results but no general homotopy classification.
\subsection{Sources}
\begin{itemize}
\item Local registry: \url{data/unsolvedmath/problems.json}, permanent UnsolvedMath ID \(6700002\), source code catalogue entry; supplied statement, background, and research summary, updated 2026-08-06.
\item Misha Gromov, \emph{101 Questions, Problems and Conjectures around Scalar Curvature} (2017), printed tag [?2], occurrence 2, PDF page 1: \url{https://www.ihes.fr/~gromov/wp-content/uploads/2018/08/101-problemsOct1-2017.pdf}. Supplied normalized PDF-text extraction accessed 2026-07-29; the source is described as incomplete.
\item The research summary mentions surgery, \(\lambda\)-invariants and families index theory without separate publications or exact result statements.
\end{itemize}
\end{document}
